% Generated draft: requires mathematical and source review.
\documentclass[11pt]{article}
\usepackage[margin=25mm]{geometry}
\usepackage{fontspec}
\setmainfont{Latin Modern Roman}
\usepackage{amsmath,amssymb,mathtools,unicode-math}
\setmathfont{Latin Modern Math}
\usepackage{enumitem,needspace,xurl,hyperref}
\setlength{\parindent}{0pt}
\setlength{\parskip}{5pt}
\setlength{\emergencystretch}{3em}
\newcommand{\N}{\mathbb{N}}
\newcommand{\Z}{\mathbb{Z}}
\newcommand{\Q}{\mathbb{Q}}
\newcommand{\R}{\mathbb{R}}
\newcommand{\C}{\mathbb{C}}
\newcommand{\F}{\mathbb{F}}
\newcommand{\A}{\mathbb{A}}
\newcommand{\PP}{\mathbb{P}}
\newcommand{\E}{\mathbb{E}}
\newcommand{\eps}{\varepsilon}
\newcommand{\dd}{\,\mathrm{d}}
\begin{document}
\section*{UnsolvedMath ID: 30006036}
\textbf{Source problem code:} catalogue entry
\subsection{Definitions and mathematical statement}
Here \(\omega\) denotes the first infinite ordinal, and ``smc \(\infty\)-category'' is read as ``symmetric monoidal \(\infty\)-category.'' For such a category \(\mathcal V\), the source denotes its scale of \(\omega\)-compact weights by \(\omega(\mathcal V)\).

The supplied conjecture is: for every symmetric monoidal \(\infty\)-category \(\mathcal V\),
\[
\mathcal V\text{ is }\omega\text{-presentable}
\quad\Longrightarrow\quad
\omega(\mathcal V)\text{ is distributive}.
\]
The expressions ``\(\omega\)-presentable,'' ``\(\omega\)-compact weight,'' ``scale,'' and ``distributive'' retain their source-specific meanings. The supplied record does not define these notions, construct \(\omega(\mathcal V)\), or state the distributivity law and its quantifiers. It also does not specify any required compatibility between the symmetric monoidal structure and presentability. Consequently, this is a schematic transcription of the conjecture, requiring the cited report's definitions and hypotheses before it becomes a self-contained mathematical statement. No particular lattice identity or colimit compatibility law is assumed here.
\subsection{Short English statement}
The conjecture asks whether every \(\omega\)-presentable symmetric monoidal \(\infty\)-category has a distributive scale of \(\omega\)-compact weights. The scale and the intended distributivity law require definitions absent from the supplied record.
\subsection{Sources}
\begin{itemize}
\item Local UnsolvedMath registry: permanent ID \(30006036\), title ``Distributivity of Compact Weight Scales,'' source problem code \texttt{catalogue entry}, in the Oberwolfach Reports Open Problems collection. The supplied proposed year is 2024; no individual proposer is identified.
\item Supplied source citation: \emph{Homotopical Algebra and Higher Structures (2025)}, Oberwolfach Reports, pp.\ 2251--2308, DOI \texttt{10.4171/owr/2024/39}. \url{https://doi.org/10.4171/owr/2024/39}. The record attributes the conjecture to this workshop report and describes its original-statement passage as retained verbatim; that description has not been independently checked.
\item Supplied literature source: ``Oberwolfach Report 39 (2024),'' \url{https://doi.org/10.4171/owr/2024/39}. Literature checked at: \texttt{2026-08-22}. This is supplied dated curation, not independently verified: it labels the problem open, describes distributivity of the scale of \(\omega\)-compact weights as an explicit conjecture, and reports that no resolution of the general statement was found. The background expressly calls for independent verification of time-sensitive status and cited claims.
\end{itemize}
\end{document}
